\documentclass[12pt, a4paper]{article}
\usepackage{amsmath, amssymb, amsthm}
\usepackage{geometry}
\geometry{a4paper, margin=1in}
\usepackage{hyperref}
\usepackage{graphicx}
\usepackage{fancyhdr}
\usepackage{cleveref}

% Theorem environments with MIT license attribution
% From: https://github.com/rollends/SE499

% Additional custom environments for your framework

% Document Information
\title{Millennium Prize Problems: A Rigorous Approach}
\author{}
\date{\today}

% Header and Footer Setup
\pagestyle{fancy}
\fancyhead[L]{Millennium Prize Problems}
\fancyhead[R]{\thepage}
\setlength{\headheight}{14.5pt} % Fix for fancyhdr warning

% Definitions and Notation Section
\newcommand{\notationsection}{
    \section*{Definitions and Notation}
    \addcontentsline{toc}{section}{Definitions and Notation}
    % Insert definitions and notations here
}

% Proof Formatting
\renewenvironment{proof}[1][Proof]{\noindent\textbf{#1.} }{\hfill $\square$\vspace{0.5em}}

\begin{document}

\maketitle

\tableofcontents
\newpage

% Introduction
\section{Introduction}
\subsection{Objective}
\textit{Provide an overview of the Millennium Prize Problems with an emphasis on CMI-required rigor.}

\subsection{Purpose}
\textit{Outline CMI's criteria for proofs, setting the foundation for the approach to each problem.}

\subsection{Definitions}
% Insert relevant definitions for the Introduction
\textit{Definitions: Theorem, Proof, CMI Standards.}

% Visual Aids (if any)
%\begin{figure}[h]
%    \centering
%    \includegraphics[width=\textwidth]{images/intro_visual.png}
%    \caption{Overview of CMI standards and Millennium Prize Problems}
%    \label{fig:intro_visual}
%\end{figure}

% Definitions and Notation
\notationsection

% 1. Birch and Swinnerton-Dyer Conjecture
\section{Birch and Swinnerton-Dyer Conjecture}
\subsection{Theorem Statement}
\begin{theorem}[Birch and Swinnerton-Dyer Conjecture]
\label{thm:bsdc}
\textit{Formal statement of the conjecture related to elliptic curves and L-functions.}
\end{theorem}

\subsection{Definitions}
% Insert relevant definitions
\textit{Definitions: Elliptic curves, L-functions, Rank of elliptic curves.}

\subsection{Proof Approach}
\textit{Apply methods in analytic number theory and computational techniques for testing the conjecture.}

\subsection{Challenges}
\textit{Analytical and computational issues in proving the conjecture formally.}

\subsection{Implications}
\textit{Advancements in number theory and cryptography could result from a proof.}

% Visual Aids (if any)
%\begin{figure}[h]
%    \centering
%    \includegraphics[width=\textwidth]{images/bsdc_visual.png}
%    \caption{Illustrates elliptic curves and L-functions}
%    \label{fig:bsdc_visual}
%\end{figure}

% 2. Hodge Conjecture
\section{Hodge Conjecture}
\subsection{Theorem Statement}
\begin{theorem}[Hodge Conjecture]
\label{thm:hodge}
\textit{Precisely state the Hodge Conjecture, emphasizing implications in algebraic geometry.}
\end{theorem}

\subsection{Definitions}
% Insert relevant definitions
\textit{Definitions: Kähler manifolds, Hodge decomposition, Algebraic cycles.}

\subsection{Proof Approach}
\textit{Detail cohomological analysis and algebraic topology techniques per CMI standards.}

\subsection{Challenges}
\textit{Connecting differential forms with algebraic cycles in a rigorous proof.}

\subsection{Implications}
\textit{Potential impacts on geometry and topology per CMI standards.}

% Visual Aids (if any)
%\begin{figure}[h]
%    \centering
%    \includegraphics[width=\textwidth]{images/hodge_visual.png}
%    \caption{Visual of cohomology and algebraic cycles}
%    \label{fig:hodge_visual}
%\end{figure}

% 3. Riemann Zeta Hypothesis
\section{Riemann Zeta Hypothesis}
\subsection{Theorem Statement}
\begin{theorem}[Riemann Hypothesis]
\label{thm:riemann}
\textit{Formalize the Riemann Hypothesis regarding non-trivial zeros of the zeta function in the critical strip.}
\end{theorem}

\subsection{Definitions}
% Insert relevant definitions
\textit{Definitions: Zeta function, Critical strip, Prime distribution.}

\subsection{Proof Approach}
\textit{Use complex analysis to understand prime distribution.}

\subsection{Challenges}
\textit{Primary obstacles related to non-trivial zeros.}

\subsection{Implications}
\textit{Impacts prime distribution if proven.}

% Visual Aids (if any)
%\begin{figure}[h]
%    \centering
%    \includegraphics[width=\textwidth]{images/rzt_visual.png}
%    \caption{Graph of the zeta function}
%    \label{fig:rzt_visual}
%\end{figure}

% 4. Navier-Stokes Equations
\section{Navier-Stokes Equations}
\subsection{Theorem Statement}
\begin{theorem}[Navier-Stokes Existence and Smoothness]
\label{thm:navier_stokes}
\textit{Statement of the Navier-Stokes existence and smoothness problem in 3D space.}
\end{theorem}

\subsection{Definitions}
% Insert relevant definitions
\textit{Definitions: Smoothness, Energy dissipation, Finite-time blowup.}

\subsection{Proof Approach}
\textit{Functional analysis and energy methods for solution smoothness.}

\subsection{Challenges}
\textit{Challenges in maintaining smoothness and energy control.}

\subsection{Implications}
\textit{Applications in fluid dynamics.}

% Visual Aids (if any)
%\begin{figure}[h]
%    \centering
%    \includegraphics[width=\textwidth]{images/ns_visual.png}
%    \caption{Energy dissipation in fluid flow}
%    \label{fig:ns_visual}
%\end{figure}

% 5. Quantum Holo-Morphic Framework
\section{Quantum Holo-Morphic Framework}
\subsection{Theorem Statement}
\begin{theorem}[Holographic Principle]
\label{thm:quantum_holo}
\textit{Present the holographic principle, especially for black hole thermodynamics.}
\end{theorem}

\subsection{Definitions}
% Insert relevant definitions
\textit{Definitions: Holography, Black hole thermodynamics, Complex time.}

\subsection{Proof Approach}
\textit{Outline principles of holographic dualities in physics per CMI standards.}

\subsection{Challenges}
\textit{Mathematical challenges in proving quantum unification principles.}

\subsection{Implications}
\textit{Implications for quantum gravity.}

% Visual Aids (if any)
%\begin{figure}[h]
%    \centering
%    \includegraphics[width=\textwidth]{images/qhmf_visual.png}
%    \caption{Illustrates holography in black hole thermodynamics}
%    \label{fig:qhmf_visual}
%\end{figure}

% 6. Yang-Mills Existence and Mass Gap
\section{Yang-Mills Existence and Mass Gap}
\subsection{Theorem Statement}
\begin{theorem}[Yang-Mills Existence and Mass Gap]
\label{thm:yang_mills}
\textit{State the Yang-Mills existence and mass gap conjecture within quantum field theory.}
\end{theorem}

\subsection{Definitions}
% Insert relevant definitions
\textit{Definitions: Mass gap, Gauge symmetry, Compactification.}

\subsection{Proof Approach}
\textit{Discuss gauge theories with a focus on precision.}

\subsection{Challenges}
\textit{Issues regarding mass gap in non-abelian gauge theories.}

\subsection{Implications}
\textit{Impact on particle physics.}

% Visual Aids (if any)
%\begin{figure}[h]
%    \centering
%    \includegraphics[width=\textwidth]{images/ymg_visual.png}
%    \caption{Illustrates gauge symmetry}
%    \label{fig:ymg_visual}
%\end{figure}

% 7. P vs NP Problem
\section{P vs NP Problem}
\subsection{Theorem Statement}
\begin{theorem}[P vs NP Problem]
\label{thm:p_vs_np}
\textit{State the P vs NP problem, focusing on computational complexity.}
\end{theorem}

\subsection{Definitions}
% Insert relevant definitions
\textit{Definitions: Computational complexity, NP-complete problems, Polynomial time.}

\subsection{Proof Approach}
\textit{Emphasize reduction techniques with rigor in class separation.}

\subsection{Challenges}
\textit{Challenges in distinguishing P and NP classes.}

\subsection{Implications}
\textit{Potential impacts on optimization and cryptography.}

% Visual Aids (if any)
%\begin{figure}[h]
%    \centering
%    \includegraphics[width=\textwidth]{images/pnp_visual.png}
%    \caption{Visual of computational complexity classes}
%    \label{fig:pnp_visual}
%\end{figure}

% Conclusion
\section{Conclusion}
\subsection{Summary of Contributions}
\textit{Summarize contributions with adherence to CMI rigor and precision.}

\subsection{Implications for Future Research}
\textit{Future research implications.}

% Visual Aids (if any)
%\begin{figure}[h]
%    \centering
%    \includegraphics[width=\textwidth]{images/conclusion_visual.png}
%    \caption{Concluding summary of document objectives}
%    \label{fig:conclusion_visual}
%\end{figure}

\end{document}
